% problem-000460 1 硅氮笼状分子合成
\section*{1. 硅氮笼状分子合成}
\textit{下列为分问。}

\subsection*{1-1}
$\mathrm { C H _ { 3 } S i C l _ { 3 } }$ 和⾦属钠在液氨中反应，得到组成为 $\mathrm { S i _ { 6 } C _ { 6 } N _ { 9 } H _ { 2 7 } }$ 的分⼦，此分⼦有⼀条三重旋转轴，所有 Si 原⼦不可区分，画出该分⼦的结构图（必须标明原⼦种类，H原⼦可不标），并写出化学反应⽅程式。

\subsection*{1-2}
⾦属钠和 $( \mathrm { C _ { 6 } H _ { 5 } ) _ { 3 } C N H _ { 2 } }$ 在液氨中反应，⽣成物中有⼀种红⾊钠盐，写出化学反应⽅程式，解释红⾊产⽣的原因。

\subsection*{1-3}
最新研究发现，⾼压下⾦属 Cs 可以形成单中⼼的 $\mathrm { C s F _ { 5 } }$ 分⼦，试根据价层电⼦对互斥理论画出 $\mathrm { C s F _ { 5 } }$ 的中⼼原⼦价电⼦对分布，并说明分⼦形状。
